\chapter{Practical Infrastructure}
\label{app:infra}

This appendix collects the software, environment and hardware details needed to
run every notebook in the course.

\section{Software stack}

\begin{center}\small
\renewcommand{\arraystretch}{1.15}
\begin{tabular}{@{}lll@{}}
\toprule
Layer & Library & Version\\
\midrule
Language & Python & 3.12\\
Deep learning & PyTorch & $\geq$2.2\\
Computer vision & torchvision & $\geq$0.17\\
Segmentation & segmentation-models-pytorch & $\geq$0.3.3\\
Metrics & torchmetrics & $\geq$1.3\\
EO datasets & torchgeo & $\geq$0.6\\
Raster I/O & rasterio & $\geq$1.3\\
Raster + xarray & rioxarray & $\geq$0.15\\
Multi-dim arrays & xarray & $\geq$2024.1\\
Vector + CRS & geopandas & $\geq$0.14\\
Augmentation & albumentations & $\geq$1.4\\
Visualisation & matplotlib / seaborn & latest\\
Interpretability & pytorch-grad-cam & $\geq$1.5\\
Notebooks & jupyterlab & $\geq$4.1\\
\bottomrule
\end{tabular}
\end{center}

\section{Environment setup}

\paragraph{uv (recommended).}
\begin{lstlisting}
curl -LsSf https://astral.sh/uv/install.sh | sh
git clone <repo-url> CNN_projects && cd CNN_projects
uv sync                 # installs from pyproject.toml into .venv
uv run jupyter lab
\end{lstlisting}

\paragraph{conda or pip.}
\begin{lstlisting}
conda env create -f environment.yml && conda activate eo-cnn && pip install -e .
# or
python3.12 -m venv .venv && source .venv/bin/activate
pip install -r requirements.txt && pip install -e .
\end{lstlisting}

\paragraph{Google Colab.} Every notebook begins with a setup cell that installs
dependencies, mounts Drive, and detects the device:
\begin{lstlisting}
IN_COLAB = 'google.colab' in str(get_ipython())
if IN_COLAB:
    !pip install -q torch torchvision torchgeo segmentation-models-pytorch \
                    torchmetrics albumentations rasterio rioxarray geopandas grad-cam tqdm
    from google.colab import drive; drive.mount('/content/drive')
    DATA_ROOT = '/content/drive/MyDrive/eo_cnn_data'
else:
    DATA_ROOT = '../data'
DEVICE = 'cuda' if torch.cuda.is_available() else 'cpu'
FAST_MODE = not torch.cuda.is_available()   # reduces epochs/data on CPU
\end{lstlisting}

\section{Repository structure}

\begin{lstlisting}
CNN_projects/
|- README.md, pyproject.toml, environment.yml, requirements.txt
|- docs/            course_design.md, infrastructure.md, capstone.md,
|                   resources.md, textbook/ (this book)
|- src/eo_cnn/      installable reference package
|   |- data/        eurosat.py, sentinel2.py, transforms.py
|   |- models/      simple_cnn.py, resnet.py, unet.py
|   |- training/    trainer.py, metrics.py
|   |- visualization/ plots.py
|- notebooks/       chapter_01_.../ ... chapter_12_.../ capstone/
|- data/            (gitignored; auto-created by torchgeo)
\end{lstlisting}

The \texttt{eo\_cnn} package is the clean reference implementation of every
concept the notebooks teach inline: \texttt{data} (EuroSAT/Sentinel-2 loaders,
tiling, albumentations pipelines), \texttt{models} (SimpleCNN, ResNet
classifiers, from-scratch and SMP U-Nets), \texttt{training} (classification and
segmentation trainers, metrics) and \texttt{visualization} (batch grids,
prediction triplets, confusion matrices, land-cover maps). The teaching
notebooks are deliberately self-contained for Colab; treat \texttt{eo\_cnn} as
the polished version to read after a chapter and to build on in the capstone.

\section{Hardware recommendations}

\begin{center}\small
\renewcommand{\arraystretch}{1.15}
\begin{tabular}{@{}llll@{}}
\toprule
Chapters & Recommended & Minimum & Colab free\\
\midrule
1--4 & CPU & CPU & CPU\\
5 & GPU 4\,GB+ & CPU (slow) & T4\\
6--8 & GPU 6\,GB+ & GPU 4\,GB & T4\\
9--10 & GPU 8\,GB+ & GPU 6\,GB & T4 (batch 4)\\
11--12 & GPU 8\,GB+ & GPU 6\,GB & T4\\
Capstone & GPU 8\,GB+ / A100 & GPU 6\,GB & Colab Pro A100\\
\bottomrule
\end{tabular}
\end{center}

\noindent The \texttt{FAST\_MODE} flag (auto-enabled without CUDA) reduces batch
size, epochs and dataset fraction so every notebook completes on CPU in under an
hour, at reduced accuracy. Dataset disk budget: EuroSAT RGB 90\,MB, EuroSAT-MS
500\,MB, LoveDA $\sim$3\,GB (all fit Colab's free disk); a BigEarthNet subset
should live on Google Drive.

\paragraph{Platform notes.} On Windows install GDAL (e.g.\ via OSGeo4W) before
\texttt{rasterio}; on Linux/Mac/Colab \texttt{pip install rasterio} works
directly. \texttt{torchgeo} caches datasets on first use and loads instantly
afterwards; all course datasets are freely accessible without registration.
